\documentclass[11pt]{article}
\usepackage[margin=1in]{geometry}
\usepackage{amsmath,amssymb}
\usepackage{booktabs}
\usepackage[hidelinks]{hyperref}
\usepackage{xcolor}
\usepackage{enumitem}

\newcommand{\NN}{\mathcal{N}}
\newcommand{\good}{\textcolor{teal}{\textbf{\checkmark}}}

\title{\bf Emergent Planning in a Relational Recurrent Core\\
\large Binding, Transition-Respecting Routing, and Planning by Value Propagation on Sokoban}
\author{Recurrent-Planner project}
\date{\today}

\begin{document}
\maketitle

\begin{abstract}
We test whether model-free RL on a \emph{state-indexed relational recurrent core} induces planning,
decomposed into a three-step chain: (1)~latent cells \textbf{bind} to environment states;
(2)~the learned content-based routing \textbf{respects the environment's transition graph} $\NN$;
and (3)~\textbf{planning emerges} over the latent states through the recurrent ``thinking'' loop.
We train an attention-LSTM core (dense attention with a sparse \mbox{1.5-entmax} normalizer, depth
$D{=}3$, variable thinking depth $d\sim\mathcal{U}\{1..6\}$) on Sokoban for 300M IMPALA steps and
probe the trained model. We find all three legs hold: (1)~task objects (box/target/agent/wall) are
linearly decodable from per-cell latents; (2)~a \emph{causal} perturbation that is blocked by a wall
influences the agent's decision $2.6\times$ less than an open one at matched pixel distance
(partial correlation $-0.32$) --- information routes \emph{around} walls, along $\NN$, not through
them; (3)~the thinking loop implements \emph{one step of generalized policy iteration}: an amortized,
graph-respecting \textbf{policy evaluation} (goal-directed value propagation) in the loop, plus a
one-step \textbf{lookahead improvement} at the actor head. The per-tick operator is a convex-average
(no-$\max$) propagation whose attention is stationary and broad, so the loop is \emph{not} value
iteration; relocating the goal reshapes the propagated latent (a $1.55\times$ shift vs.\ a matched
transition change), so it is \emph{not} a reward-agnostic successor representation either --- it
evaluates a goal-directed value field whose converged form is a recursive Bellman fixed point. The
policy then agrees with $\arg\max_a[r_a+\gamma V(s'_a)]$ over its own value above chance ($0.41$ vs.\
$0.25$), rising with thinking depth ($0.34\!\to\!0.41$); the horizon is amortized (no multi-step
internal trajectory is stored). Critically, the relational structure of (1)--(2) \emph{emerges from RL
with no architectural mask}.
\end{abstract}

\section{Setup}

\paragraph{Core.} The recurrent backbone is a state-indexed attention-LSTM: each latent cell is the
$C{=}32$-channel vector at a spatial board square $s$ (NHWC carry over $S{=}H{\times}W$ tokens). Per
``thinking tick'' every cell updates by attending over the board (4 heads), aggregating neighbour
values, and applying an LSTM gate. We use \textbf{dense} attention (no locality mask) with a
\textbf{sparse 1.5-entmax} normalizer, depth $D{=}3$ stacked cells (\texttt{n\_recurrent}), and up to
$K{=}6$ ticks. The map to the relational-update formalism is
$h^{k+1}(s')=\mathrm{AGG}_{(s,s')}\,F_\phi(h^k(s))$, with $\NN$ the attention support, $F_\phi$ the
value projection, and the gate the LSTM update.

\paragraph{Three training requirements.} The investigation rests on three mechanisms, all verified:
\begin{enumerate}[leftmargin=1.5em,itemsep=1pt]
  \item \textbf{Random thinking depth.} A depth $d\sim\mathcal{U}\{1..6\}$ is sampled per rollout; the
        learner replays the recurrence at the same $d$ (V-trace consistency).
  \item \textbf{$k$-step recurrent backprop.} A gated scan runs the static $K$ ticks but freezes the
        carry past $d$, so gradients flow through exactly $d$ weight-tied applications. Verified
        \emph{bit-exact}: $\nabla(\text{gated }d) = \nabla(\text{native $d$-tick model})$ to $0.0$,
        and the recurrent-cell gradient is exactly $0$ at $d{=}0$ and grows monotonically with $d$.
  \item \textbf{Sparsity via entmax (architectural).} Attention weights are 1.5-entmax (true zeros),
        not a soft entropy penalty --- one module, no loss term, no scan plumbing.
\end{enumerate}

\paragraph{Training outcome.} On unfiltered-train Sokoban (IMPALA, 300M steps) the model learns
cleanly: online train solve rate $\approx\!83\%$, episode length $76\!\to\!40$, value
$\text{var\_explained}\approx0.99$, policy entropy $1.37\!\to\!0.29$. On \texttt{valid\_medium}
(generalization) the greedy full-depth solve rate is $\approx\!31\%$. A faithful sweep of the trained
thinking depth $d$ (Table~\ref{tab:think}) shows solve rate \emph{rises} with thinking, confirming the
loop does useful computation.

\begin{table}[h]\centering
\caption{Thinking-curve: \texttt{valid\_medium} solve rate vs.\ inner thinking depth $d$ (faithful
$n_{\text{active}}$ sweep; whole episode run at depth $d$). Thinking helps; mild interior peak at $d{=}4$.}
\label{tab:think}
\begin{tabular}{lccccccc}
\toprule
depth $d$ & 0 & 1 & 2 & 3 & 4 & 5 & 6 \\
\midrule
solve (160M ckpt) & 0.000 & 0.258 & 0.314 & 0.325 & \textbf{0.336} & 0.329 & 0.319 \\
solve (140M ckpt) & 0.000 & 0.223 & 0.288 & 0.310 & 0.316 & --- & 0.319 \\
\bottomrule
\end{tabular}
\end{table}

\section{Step 1: Latent states bind to environment states \good}

We linearly decode the tile type at square $s$ from its latent $h(s)$ (per-object balanced accuracy,
chance $=0.5$). All task-relevant objects are decodable well above chance, including the rare ones
(box/target/agent are $1$--$4\%$ of squares). Binding is established immediately (flat over ticks),
consistent with it being read from the input embedding rather than constructed by the loop.

\begin{table}[h]\centering
\caption{Per-object binding: balanced accuracy of a linear probe $h(s)\to$ is-object (chance $0.5$).}
\label{tab:bind}
\begin{tabular}{lcccc}
\toprule
object & wall & box & target & agent \\
base rate & 0.69 & 0.04 & 0.04 & 0.01 \\
\midrule
balanced acc & 0.76--0.78 & 0.76 & 0.70 & 0.74--0.76 \\
\bottomrule
\end{tabular}
\end{table}

\noindent\emph{(Note: an earlier floor-dominated 5-way aggregate probe sat at chance and was
uninformative; the per-object balanced probe is the correct measurement.)}

\section{Step 2: Routing respects the transition graph $\NN$ \good}

\paragraph{Recovery of $\NN$ (correlational).} With no locality mask, the attention nonetheless
concentrates its \emph{mass} on the one-step move-reachable (wall-masked) neighbours far above chance
(Table~\ref{tab:recoverN}): $7$--$13\times$, vs.\ $3.1\times$ for the earlier dense-softmax core. The
entmax support is broad ($\sim$44--88 nonzero keys), so this is mass-concentration with a long tail,
not hard sparsity.

\begin{table}[h]\centering
\caption{Recovery of $\NN$: attention mass on wall-masked one-step neighbours vs.\ uniform chance
($0.027$), per stacked cell (final tick, 160M).}
\label{tab:recoverN}
\begin{tabular}{lccc}
\toprule
 & cell 0 & cell 1 & cell 2 \\
\midrule
mass on reachable & 0.351 & 0.320 & 0.189 \\
lift over chance & $13.0\times$ & $11.8\times$ & $7.0\times$ \\
\bottomrule
\end{tabular}
\end{table}

\paragraph{Through-walls (causal).} Mass-on-$\NN$ alone does not prove the \emph{computation} respects
transitions --- a global tail could carry information through walls. We test it causally: flip one
floor cell to a wall at Euclidean distance $\in[3.5,7]$ from the agent (beyond the conv receptive
field, so the agent's embedding is unaffected) and measure the agent-cell latent shift, comparing
cells \emph{blocked} by walls (graph $\gg$ Euclidean / unreachable) vs.\ \emph{open} ones at matched
pixel distance (Table~\ref{tab:walls}). A perturbation behind a wall moves the agent's latent only
$0.38\times$ as much (partial correlation of influence with graph-distance, controlling Euclidean,
$=-0.32$). \textbf{Information routes around walls, along $\NN$, not through Euclidean space.}

\begin{table}[h]\centering
\caption{Through-walls causal test (300M; matched Euclidean distance $\approx4.5$). Influence $=$
final agent-cell $\lVert\Delta h\rVert$ from the perturbation.}
\label{tab:walls}
\begin{tabular}{lccc}
\toprule
 & influence $\lVert\Delta h\rVert$ & onset (tick) & Euclidean dist. \\
\midrule
open (navigable path) & 0.857 & 2.29 & 3.9--4.6 \\
blocked by wall       & 0.323 & 1.73 & 4.5--5.4 \\
\midrule
\multicolumn{4}{l}{blocked/open ratio $=0.38$; \quad partial corr$(\,\text{infl},\text{graph}\mid\text{Euclid})=-0.32$} \\
\bottomrule
\end{tabular}
\end{table}

\noindent A separate perturbation sweep shows the influence \emph{onset} (timing) is roughly
distance-independent ($\approx 2$ ticks): information arrives fast/globally, but \emph{how much}
arrives is gated by graph connectivity. The routing is thus fast (parallel-like) yet
transition-respecting in magnitude.

\section{Step 3: Planning as two-step policy iteration \good}

The thinking loop realizes \emph{one step of generalized policy iteration, split across the
architecture}: (3a)~the recurrent loop performs amortized, graph-respecting \textbf{policy evaluation}
(goal-directed value propagation along $\NN$), and (3b)~the actor head performs a one-step
\textbf{lookahead policy improvement}. The loop is \emph{not} value iteration and \emph{not}
successor-representation construction; we pin each below.

\subsection*{(3a) The loop evaluates: goal-directed value propagation}

\paragraph{One tick is a damped propagation, not a Bellman max.} The trained core uses
\texttt{readout="softmax"}, \texttt{attn\_norm="entmax15"}, \texttt{directional\_value=False}: the
per-tick aggregation $z=(A\,v)\,W_{\!out}$ is a \emph{convex average} $\mathbb{E}_{s'\sim A}[\cdot]$,
\emph{not} a $\max$ (the codebase's soft-Bellman \texttt{maxplus} operator is off for this checkpoint).
Linearizing the value channel, one weight-tied tick and its fixed point are
\begin{equation}
  c \;\leftarrow\; \gamma_{\text{eff}}\,A\,c + r_{\text{eff}}
  \qquad\Longrightarrow\qquad
  c^\star=(I-\gamma_{\text{eff}}A)^{-1}r_{\text{eff}}=\textstyle\sum_{i\ge0}(\gamma_{\text{eff}}A)^i r_{\text{eff}},
  \label{eq:resolvent}
\end{equation}
the discounted \textbf{policy-evaluation resolvent} (equivalently the successor-representation kernel):
$A\leftrightarrow\gamma P^\pi$ (the entmax routing on $\NN$), the per-channel forget gate
$\sigma(g_f)\leftrightarrow$ the discount, and the $[x;\text{prev}]W_{\!in}$ injection
$\leftrightarrow$ the reward source $r$. No $\max_a$ appears in the loop.

\paragraph{The attention operator is stationary (E1).} Value iteration would sharpen $A$ toward the
greedy successor as the estimate improves. It does not: across ticks $A$ stays broad and stationary
(Table~\ref{tab:e1}), consistent with a \emph{fixed} transition operator $\gamma P^\pi$ rather than a
sharpening optimality backup.

\paragraph{Policy evaluation, not a successor representation (E2).} The two members of
Eq.~\eqref{eq:resolvent} differ in \emph{where reward enters}: policy-evaluation bakes the goal into
the iterate $c^\star$; a successor representation stores reward-agnostic occupancy and injects reward
only at the critic. We relocate the target to a far floor cell (move the reward, keep $P$) and measure
the recurrent latent shift outside the conv receptive field, against a matched far wall-flip (change
$P$, keep reward). Moving the goal shifts the propagated latent $h_2$ \emph{more} than a transition
change does (Table~\ref{tab:e2}: far-field ratio $1.55$, agent-cell $1.34$), with a large readout-value
shift. A successor representation predicts a $\approx0$ far-field shift; this falsifies it.

\paragraph{Genuine but amortized iteration.} The hidden contracts geometrically at $\rho\approx0.89$
(effective horizon $\sim$9 ticks) and is only $\sim$60\% converged at the trained $K{=}6$: the loop does
real inference-time iteration, but its effective horizon is far shorter than the task's $\gamma{=}0.97$
horizon, so the long-range solution is \emph{amortized} into the weights and the loop runs a few
refinement sweeps. The propagation is moreover \emph{parallel}: $A$ is broad ($\sim$85/100 keys) and (\S2)
a perturbation's influence onset is roughly distance-independent ($\sim$2--3 ticks) with graph-gated
\emph{magnitude} --- a parallel graph-weighted spread, not a serial one-ring-per-tick wavefront.

\paragraph{The converged value is a recursive Bellman fixed point.} Along the greedy trajectory the
optimality residual $|V(s_k)-\max_a[r_a+\gamma V(s'_{k,a})]|$ is low and \emph{flat} across depth
(Table~\ref{tab:bellman}), with a small constant \emph{negative} (soft/entropy-regularized) gap: $V$
approximately satisfies the recursive Bellman equation everywhere. This characterizes the
\emph{converged value} (depth-agnostic); it does not by itself identify the per-tick operator --- E1/E2 do.

\begin{table}[h]\centering
\caption{E1 (300M): the per-tick operator is stationary and broad and contracts slowly --- not a
sharpening Bellman backup. Ticks $1\!\to\!12$ (trained $K{=}6$).}
\label{tab:e1}
\begin{tabular}{lcc}
\toprule
quantity & tick 1 & tick 12 \\
\midrule
attention top-1 mass / query        & 0.156 & 0.159 \\
entmax support (\#keys $>10^{-6}$)  & 87 & 85 \\
operator drift $\cos(A_t,A_{t-1})$  & 0.95 & 0.995 \\
relative update $\lVert\Delta h_t\rVert/\lVert h_t\rVert$ & 0.48 & 0.10 \\
\bottomrule
\end{tabular}\\[3pt]
{\footnotesize Contraction $\rho\approx0.89$ (effective horizon $\sim$9 ticks); $\sim$60\% converged at $K{=}6$.}
\end{table}

\begin{table}[h]\centering
\caption{E2 (300M; 230 verified relabels): relocating the goal moves the propagated latent \emph{more}
than a matched transition change --- the goal is in the iterate (policy-evaluation), not a
reward-agnostic successor representation (which predicts $\approx0$).}
\label{tab:e2}
\begin{tabular}{lccc}
\toprule
relative latent shift $\lVert\Delta h_2\rVert/\lVert h_2\rVert$ & reward-move & wall-flip & ratio \\
\midrule
far field (outside conv RF of both edits) & 0.349 & 0.225 & 1.55 \\
agent cell (clean of both edits)          & 0.273 & 0.204 & 1.34 \\
\midrule
readout value shift $|\Delta V|/\mathrm{std}(V)$ & 0.84 & 0.45 & --- \\
\bottomrule
\end{tabular}
\end{table}

\begin{table}[h]\centering
\caption{Bellman self-consistency (300M): optimality residual (normalized by $\mathrm{std}(V)$) along the
greedy trajectory --- low and flat $\Rightarrow$ a recursive (multi-step) fixed point.}
\label{tab:bellman}
\begin{tabular}{lcccc}
\toprule
depth $k$ & 0 & 1 & 2 & 3 \\
\midrule
$|V-\max_a(r+\gamma V(s'))|/\mathrm{std}(V)$ & 0.19 & 0.25 & 0.25 & 0.22 \\
\bottomrule
\end{tabular}
\end{table}

\subsection*{(3b) The head improves: one-step lookahead}

\paragraph{Lookahead consistency (the decisive test).} Planning means evaluating actions at the
states they transition into. For each board we step into each action's successor $s'_a$, read the
model's \emph{own} value $V(s'_a)$, and ask whether the chosen action equals the one-step lookahead
$a^\star \stackrel{?}{=} \arg\max_a\,[\,r_a+\gamma V(s'_a)\,]$. The policy is consistent with this
lookahead well above chance, and --- crucially --- \textbf{consistency rises with thinking depth}
(Table~\ref{tab:lookahead}): extra ticks move the action toward the value-maximizing successor
($13.7\%$ of boards have their action flipped to the lookahead choice by the loop). This is the
mechanistic correlate of the rising thinking-curve.

\begin{table}[h]\centering
\caption{Lookahead consistency (300M): agreement of the policy with one-step lookahead over its own
value. Chance $=1/n_{\text{act}}=0.25$.}
\label{tab:lookahead}
\begin{tabular}{lccccccc}
\toprule
 & full & \multicolumn{6}{c}{by thinking depth $d$}\\
\cmidrule(l){3-8}
 & depth & $d{=}1$ & 2 & 3 & 4 & 5 & 6 \\
\midrule
agreement with $\arg\max_a[r_a+\gamma V(s'_a)]$ & 0.41 & 0.34 & 0.38 & 0.39 & 0.38 & 0.42 & 0.41 \\
\bottomrule
\end{tabular}
\end{table}

\paragraph{No stored multi-step trajectory.} Decoding the model's \emph{executed} future actions
from the policy readout, only the immediate action $a_0$ is represented (and it sharpens over ticks,
$0.70\!\to\!0.97$); the future plan $a_1,\dots,a_5$ sits at chance. So there is \emph{no multi-step
internal trajectory} being built or optimized --- the long horizon is amortized into the value
function, and the loop refines a one-step, value-based decision (a policy-improvement primitive),
not a multi-step forward search.

\section{Methodology and caveats}

\paragraph{Faithful recompute.} Because the model's own forward (the offset-tied \texttt{rel\_bias}
gather under \texttt{nn.scan}) triggers a tracer error when traced/jitted, all probes recompute the
$D{=}3$ entmax stack in plain JAX from the loaded parameters and read out logits/value with the
trained head matmuls. The recompute is validated to reproduce the model's own final hidden state to
$1.8\times10^{-7}$.

\paragraph{Corrected probes.} Two early measurements were methodologically wrong and were corrected;
we report this for transparency. (i)~A flat $0$-vs-$8$ ``steps-to-think'' eval was \emph{not} the
thinking axis --- the faithful $n_{\text{active}}$ sweep (Table~\ref{tab:think}) is. (ii)~Decoding the
future plan from the \emph{initial} state's representation cannot detect planning, because future
actions are taken at future states; the successor-value lookahead test (Table~\ref{tab:lookahead}) is
the correct test, and it flips the conclusion from ``reactive'' to ``value-based lookahead.''

\paragraph{Limits.} Probes are linear (a nonlinear probe might recover more structure); the causal
tests are on the final checkpoint (an emergence-over-training sweep is the natural next figure);
lookahead is measured at one step (the value amortizes deeper horizons); ``respects walls'' has a
mild confound (walled cells may be intrinsically less plan-relevant, though that is itself
graph-aware computation).

\section{Conclusion}

Model-free RL on a relational recurrent core yields, \emph{with no imposed locality mask}:
\textbf{(1)}~latent cells that bind to environment states; \textbf{(2)}~content-based routing whose
information flow respects the transition graph $\NN$ (causally: it routes around walls); and
\textbf{(3)}~a thinking loop that realizes one step of generalized policy iteration --- amortized,
graph-respecting \emph{policy evaluation} (goal-directed value propagation) in the loop, and a one-step
\emph{lookahead improvement} at the actor head. The structure of (1)--(2) is \emph{emergent}, which is
the central claim --- masking the attention would hard-code $\NN$ and make the result circular. The
planning (3) is genuine but \emph{shallow and amortized}: the loop propagates a goal-directed value
field through a fixed, broad transition operator (a few parallel refinement sweeps; effective horizon
$\sim$9 ticks), \emph{not} deep multi-step search, trajectory optimization, or value iteration. Inducing
deeper, explicitly multi-step planning (e.g.\ via depth curricula or convergence-promoting
``deep-thinking'' objectives) is the natural next direction.

\end{document}
